\documentclass[a4paper, 12pt]{article}

\usepackage[portuges]{babel}
\usepackage[utf8]{inputenc}
\usepackage{amsmath}
\usepackage{indentfirst}
\usepackage{graphicx}
\setlength{\marginparwidth}{2cm}
\usepackage[colorinlistoftodos]{todonotes}
\usepackage{verbatim}
\usepackage{textcomp}
\usepackage{gensymb}
\usepackage{tikz}
\usetikzlibrary{arrows.meta,calc}

\begin{document}

\begin{titlepage}
\begin{center}
\textbf{\LARGE Universidade de Brasília}\\[0.5cm] 
\textbf{\large Departamento de Ciência da Computação}\\[0.2cm]
\vspace{20pt}
\includegraphics{Logo_UnB.png}\\[1cm]

\par
\vspace{32pt}
\textbf{\LARGE Projeto Final de Sistemas de Informação}\\
\vspace{30pt}
\textbf {\Large Autores:}\\[0.2cm]
\Large {Davi Galvão Guerra		24/1038577}\\[0.1cm]
\Large {Élvis Corrêa Miranda Júnior	24/1038700}\\[0.1cm]
\Large {Gustavo Alves Dias	24/1020779}\\[0.1cm]
\Large {Leonardo Mileo Moreira Krauss Guimarães    24/1032690}\\[0.1cm]
\Large {Pedro Marcinoni Varão Ribeiro Viera da Rosa    24/1002396}\\[0.1cm]

\end{center}

\par
\vfill
\begin{center}
{{\normalsize Brasília}\\
{\normalsize \today}}
\end{center}
\end{titlepage}

%Sumário
\newpage
\tableofcontents
\thispagestyle{empty}
%End Sumário

\newpage
\section{Resumo}

Este trabalho aplica o processo de Engenharia de Requisitos ao problema do \textit{Monitoramento Multimodal de Cultivos}, no qual um robô autônomo percorre uma área agrícola para coletar dados de um conjunto predefinido de alvos georreferenciados, com apoio de um drone que fornece informações sobre a topologia do local. O contexto envolve restrições relevantes: verificação de bateria antes e durante a missão, tratamento de alvos inacessíveis por meio de posições alternativas, avaliação de qualidade dos dados com até duas novas tentativas, retorno seguro à base com retomada da missão após a recarga, visita a alvos opcionais apenas quando houver tempo e energia, e possibilidade de monitoramento e intervenção humana.

O desenvolvimento seguiu as fases de Iniciação, Investigação, Decisão, Formulação e Validação. Foram utilizados Diagrama de Contexto e \textit{Stakeholder Onion Diagram}, Cenários de Domínio, Mapa de Impacto, Especificação por Exemplos (\textit{Example Map} com BDDs) e \textit{Goal Models} em KAOS, evoluídos nas versões v1, v2 e final, com uso de \textit{goal patterns}. O resultado é um conjunto de objetivos, requisitos de máquina e suposições de domínio claros, relevantes, completos e rastreáveis, discutidos com as partes interessadas em \textit{workshop} de validação.

\textbf{Palavras-chave:} Engenharia de Requisitos, Goal Models, Impact Map, Behaviour Driven Development.

\section{Introdução}

O presente relatório descreve o processo de Engenharia de Requisitos aplicado ao problema do \textit{Monitoramento Multimodal de Cultivos}. Na agricultura de precisão, decisões como irrigação, adubação e controle de pragas dependem de dados frequentes, confiáveis e georreferenciados sobre a lavoura. A coleta manual desses dados é lenta, cara e pouco repetível, o que motiva o uso de um robô autônomo que percorre a área cultivada e coleta dados em um conjunto predefinido de alvos georreferenciados, com apoio de um drone que fornece informações sobre a topologia do local.

A complexidade do problema reside nas restrições operacionais do ambiente agrícola, com destaque para o gerenciamento autônomo de energia. A bateria do robô deve ser verificada antes do início da missão e também durante a execução, de modo que sempre reste carga para o retorno seguro à base. Ao chegar à base, o robô recarrega e retoma a missão, sem repetir alvos já concluídos. Além disso, um alvo pode estar inacessível por obstáculo, terreno intransitável ou alagamento. Nesse caso, o sistema deve recorrer a posições alternativas de coleta em vez de abandonar o alvo imediatamente.

Há ainda requisitos de qualidade e de adaptação. Cada coleta precisa ter a qualidade avaliada, e uma leitura insatisfatória pode ser repetida até duas vezes antes de ser registrada como falha. Os alvos opcionais só devem ser visitados quando sobrarem tempo e energia, sem comprometer os alvos obrigatórios nem o retorno à base. Por fim, o sistema deve permitir que um operador humano acompanhe a missão e intervenha quando necessário.

Essas restrições interagem entre si: visitar um alvo opcional consome a energia que poderia ser necessária para repetir uma coleta, e uma posição alternativa mais distante reduz a margem de retorno. Para traduzir essas regras em especificações claras e verificáveis, este documento apresenta a modelagem do sistema ao longo das fases de Iniciação, Investigação, Decisão, Formulação e Validação. A abordagem parte do entendimento das partes interessadas e das fronteiras do sistema, passa pela investigação de cenários e pela definição de objetivos, e chega à especificação por exemplos e à validação com as partes interessadas.

\section{Fundamentação Teórica}

Nesta seção, apresenta-se a base teórica das técnicas e modelos adotados ao longo das fases do projeto:

\begin{itemize}
    \item \textbf{Fase 1 (Iniciação):} utilização do \textit{Context Diagram} para mapear a fronteira do sistema e suas interações com os atores, acompanhado do \textit{Stakeholder Onion Diagram} para classificar as partes interessadas em camadas de envolvimento.
    \item \textbf{Fase 2 (Investigação):} emprego de \textit{Domain Scenarios} para narrar o fluxo de eventos e o tratamento de exceções, e construção da primeira versão do \textit{Goal Model} em KAOS para traduzir as regras do domínio em objetivos.
    \item \textbf{Fase 3 (Decisão):} aplicação do \textit{Impact Map} para alinhar os objetivos com o comportamento dos atores e as funcionalidades do sistema, resultando na evolução para o \textit{Goal Model} (v2), com uso de \textit{goal patterns}.
    \item \textbf{Fase 4 (Formulação):} uso de \textit{Specification by Example} (\textit{Example Maps} e BDDs) para criar critérios de aceite executáveis, consolidando os requisitos no \textit{Goal Model} final.
    \item \textbf{Fase 5 (Validação):} discussão dos resultados com as partes interessadas em \textit{workshop}.
\end{itemize}

\section{Fase 1: Iniciação}

\subsection{Diagrama de Contexto}

O Diagrama de Contexto delimita a fronteira do Sistema de Controle do Robô de Monitoramento Multimodal de Cultivos (RMMC) e mostra os fluxos trocados com os atores e sistemas externos. O RMMC é tratado como uma única caixa, sem detalhar sua estrutura interna.

Dois atores humanos interagem com o sistema. O \textbf{Planejador da Missão} envia os \textit{dados da missão} (alvos georreferenciados obrigatórios e opcionais). O \textbf{Controlador da Missão} recebe \textit{alertas de bateria e status da missão} e, quando necessário, devolve \textit{ajustes de alvos}, \textit{comandos de recarga} e \textit{comandos de controle}.

Dentro da fronteira estão ainda dois sistemas: o \textbf{Drone}, que fornece os \textit{dados de topologia} do local, e a \textbf{Base de Recarga}, que fornece a \textit{corrente de recarga} ao robô. Fora dela, o \textbf{Software de Análise (\textit{Dashboard})} recebe os \textit{dados coletados} nos alvos.

A Figura~\ref{fig:contexto} apresenta o diagrama. A linha tracejada delimita a fronteira do sistema, que contém o RMMC, o Controlador da Missão, o Drone e a Base de Recarga; o Planejador da Missão e o Software de Análise ficam fora dela. Os bonecos de palito representam atores humanos e as caixas representam sistemas.

\begin{figure}[htbp]
\centering
\resizebox{\textwidth}{!}{%
\begin{tikzpicture}[
    >=Stealth, font=\small,
    sys/.style={draw, thick, minimum width=2.4cm, minimum height=1.2cm, align=center},
    ext/.style={draw, thick, minimum height=1cm, align=center},
    flow/.style={->, thick},
    lbl/.style={font=\scriptsize, align=center, fill=white, inner sep=1pt}
]
% fronteira do sistema
\draw[dashed, thick] (1.5,-0.5) ellipse (6.2cm and 4.8cm);
\node[sys] (rmmc) at (0,0) {\textbf{RMMC}};

% atores (bonecos de palito)
\newcommand{\actor}[3]{%
  \begin{scope}[shift={(#2)}, thick]
    \draw (0,0.55) circle (0.2);
    \draw (0,0.35)--(0,-0.15);
    \draw (-0.3,0.2)--(0.3,0.2);
    \draw (0,-0.15)--(-0.25,-0.6);
    \draw (0,-0.15)--(0.25,-0.6);
  \end{scope}
  \node[minimum width=1cm, minimum height=1.4cm] (#1) at (#2) {};
  \node[align=center, above=0.05cm of #1] {#3};}
\actor{pm}{-8.5,2.5}{Planejador\\da Missão}
\actor{cm}{2.2,2.0}{Controlador\\da Missão}

% sistemas externos
\node[ext] (dr) at (5.2,-2.4) {Drone};
\node[ext] (br) at (1,-3.6) {Base de\\Recarga};
\node[ext] (sa) at (-7.5,-4.2) {Software de Análise\\(\textit{Dashboard})};

% fluxos
\draw[flow] (pm) -- node[lbl, sloped, above] {Dados da missão} (rmmc);
\draw[flow] (rmmc) to[bend left=20] node[lbl, above left, pos=0.4] {Alertas de bateria\\e status da missão} (cm);
\draw[flow] (cm) to[bend left=20] node[lbl, below right, pos=0.5] {Ajuste de alvos /\\comandos de recarga /\\comandos de controle} (rmmc);
\draw[flow] (dr) -- node[lbl, sloped, above] {Dados de topologia} (rmmc);
\draw[flow] (br) -- node[lbl, right] {Corrente\\de recarga} (rmmc);
\draw[flow] (rmmc) -- node[lbl, sloped, above] {Dados coletados} (sa);
\end{tikzpicture}}
\caption{Diagrama de Contexto do RMMC.}
\label{fig:contexto}
\end{figure}


\subsection{Stakeholder Onion Diagram}

O \textit{Stakeholder Onion Diagram} (Figura~\ref{fig:onion}) classifica as partes interessadas em camadas concêntricas de envolvimento com o RMMC: \textit{Operational Area} (usuários diretos e tecnologias de interface), \textit{Containing Business} (cliente e beneficiários funcionais) e \textit{Wider Environment} (afetados, reguladores e fornecedor). A \textit{Development Team} aparece à parte.

\begin{figure}[htbp]
\centering
\resizebox{!}{14cm}{%
\begin{tikzpicture}[
    font=\scriptsize,
    tech/.style={draw, thick, fill=yellow!35, align=center, minimum width=2.4cm, minimum height=0.9cm},
    org/.style={draw, thick, align=center, text width=1.8cm, minimum width=1.9cm, minimum height=0.9cm},
    ring/.style={font=\scriptsize\bfseries, align=center}
]
\newcommand{\pessoa}[3]{%
  \begin{scope}[shift={(#2)}, thick]
    \draw (0,0.55) circle (0.2);
    \draw (0,0.35)--(0,-0.15);
    \draw (-0.3,0.2)--(0.3,0.2);
    \draw (0,-0.15)--(-0.25,-0.6);
    \draw (0,-0.15)--(0.25,-0.6);
  \end{scope}
  \node[align=center, anchor=north, text width=2.3cm] at ($(#2)+(0,-0.65)$) {#3};}

% camadas
\draw[thick, fill=gray!8]  (0,0) circle (8.6cm);
\draw[thick, fill=gray!16] (0,0) circle (6.6cm);
\draw[thick, fill=gray!26] (0,0) circle (4.3cm);
\draw[thick, fill=white]   (0,0) circle (1.7cm);

\node[align=center, text width=2.8cm] at (0,0) {\textbf{Sistema de Controle do Robô de Monitoramento Multimodal de Cultivos} (RMMC)};

% Operational Area
\node[ring] at (90:3.6) {OPERATIONAL AREA};
\pessoa{c1}{150:3.0}{Controlador\\da Missão\\\textit{(user)}}
\pessoa{c2}{30:3.0}{Planejador\\da Missão\\\textit{(user)}}
\node[tech] at (210:3.0) {\textbf{Drone}\\interface de informação\\(topologia do local)};
\node[tech] at (330:3.0) {\textbf{Software de Análise}\\interface de informação\\(\textit{dashboard} da missão)};
\node[tech] at (270:3.2) {\textbf{Base de Recarga}\\interface de energia\\(recarga e retorno seguro)};
\node[font=\scriptsize\itshape] at (270:2.2) {Interfacing Technologies};

% Containing Business
\node[ring] at (90:5.8) {CONTAINING BUSINESS};
\pessoa{b1}{180:5.45}{Dono da Área Agrícola\\(Produtor Rural /\\Empresa Agrícola)\\\textit{(client / customer)}}
\pessoa{b2}{0:5.45}{Funcionários da Fazenda\\(Eng. Agrônomo e\\Analista de Dados)\\\textit{(functional beneficiaries)}}

% Wider Environment
\node[ring] at (90:7.7) {WIDER ENVIRONMENT};
\pessoa{w1}{235:7.3}{Equipe de Manutenção Física\\da Fazenda\\\textit{(negatively affected)}}
\node[org, draw=red!70!black, fill=red!10] at (25:7.6) {Órgãos Reguladores\\(ANATEL / ANAC)\\\textit{(regulators)}};
\node[org, draw=blue!70!black, fill=blue!10] at (155:7.6) {Fabricante das peças\\do Robô\\\textit{(fornecedor de hardware)}};

% Development Team
\node[ring] at (7.3,-6.2) {DEVELOPMENT TEAM};
\pessoa{d1}{7.3,-7.2}{Grupo de Alunos\\(desenvolvedores do robô)\\\textit{(developers)}}
\end{tikzpicture}}
\caption{\textit{Stakeholder Onion Diagram} do RMMC.}
\label{fig:onion}
\end{figure}

\section{Fase 2: Investigação}

\subsection{Domain Scenarios}

\subsection{Goal Model (v1)}

\section{Fase 3: Decisão}

\subsection{Impact Map}

\subsection{Goal Model (v2) e Goal-Feature Mapping}

\section{Fase 4: Formulação}

\subsection{Specification by Example: Example Mapping}

\subsection{Specification by Example: BDD (Gherkin)}

\subsection{Goal Model (final)}

\section{Fase 5: Validação}

\section{Conclusão}


\section{Referência}

\noindent

\end{document}